% 雨果·斯坦豪斯（Hugo Steinhaus）（综述）
% license CCBYSA3
% type Wiki

本文根据 CC-BY-SA 协议转载翻译自维基百科\href{https://en.wikipedia.org/wiki/Hugo_Steinhaus}{相关文章}。

\begin{figure}[ht]
\centering
\includegraphics[width=6cm]{./figures/2db84a4d97edc48c.png}
\caption{} \label{fig_YGstgs_1}
\end{figure}
雨果·迪奥尼济·斯坦豪斯（Hugo Dyonizy Steinhaus，英语发音：/ˈhjuːɡoʊ ˈstaɪnˌhaʊs/，波兰语发音：[ˈxuɡɔ ˈʂtajnxaws]，1887年1月14日－1972年2月25日）是一位波兰数学家和教育家。斯坦豪斯于1911年在哥廷根大学师从大卫·希尔伯特（David Hilbert）获得博士学位，后来成为利沃夫扬·卡齐米日大学（现今乌克兰利维夫）的教授，并在那里帮助建立了后来被称为“利沃夫数学学派”（Lwów School of Mathematics）的学术中心。他被誉为“发现”数学家斯特凡·巴拿赫（Stefan Banach）的人，与巴拿赫共同在泛函分析领域作出了重要贡献，特别是提出了著名的巴拿赫–斯坦豪斯定理（Banach–Steinhaus theorem）。

第二次世界大战后，斯坦豪斯在弗罗茨瓦夫大学（Wrocław University）数学系的建立以及战后波兰数学界的重建中发挥了重要作用。斯坦豪斯一生发表了约170篇科学论文与著作，他在数学的多个分支领域留下了深远影响，包括泛函分析、几何学、数理逻辑以及三角学。值得注意的是，他被认为是博弈论（game theory）和概率论（probability theory）的早期奠基者之一，为后续学者发展更系统的理论框架奠定了基础。
\subsection{早年生活与求学经历}
斯坦豪斯于1887年1月14日出生于奥匈帝国的亚斯沃（Jasło）一个有犹太血统的家庭。\(^\text{[1]}\)他的父亲博古斯瓦夫（Bogusław）是一名地方实业家，拥有一家砖厂，同时也是一位商人；他的母亲名为埃韦丽娜（Ewelina），娘家姓利普希茨（Lipschitz）。斯坦豪斯的叔叔伊格纳奇·斯坦豪斯（Ignacy Steinhaus）\(^\text{[2]}\)是一名律师，也是“波兰议会俱乐部”（Koło Polskie）的成员，并担任加利西亚和洛多梅里亚王国（Kingdom of Galicia and Lodomeria）地区议会——加利西亚议会（Galician Diet）的代表。
